\chapter{Axioms}
\label{chap:axioms}

\section{Minimal foundation}

The framework rests on four claims, the first three of which are essentially
established physics and the fourth of which is a minimal natural extension.

\begin{axiom}[Quantum mechanics applies]\label{ax:qm}
The supraverse is described by quantum field theory. Vacuum states have nonzero
fluctuations governed by the uncertainty principle. Over infinite spacetime,
fluctuations of arbitrary magnitude occur with probability one.
\end{axiom}

\begin{axiom}[Special relativity applies]\label{ax:sr}
Spacetime has Lorentzian signature. Causal structure is determined by light
cones. CPT symmetry holds for any local Lorentz-invariant QFT.
\end{axiom}

\begin{axiom}[Einstein--Cartan gravity applies]\label{ax:ec}
The connection on spacetime is allowed to be asymmetric. The antisymmetric part
(torsion) couples to fermion spin via the axial current. In the absence of
fermionic matter, the theory reduces to standard general relativity. With
fermions, torsion produces an effective interaction between spins that scales as
the square of the spin density, and is therefore always positive (repulsive at
high spin density) regardless of relative spin orientation.
\end{axiom}

\noindent The Einstein--Cartan Lagrangian, in shorthand:
\[
  \mathcal{L}_{EC} = \frac{1}{2\kappa}(R - 2\Lambda)
  + \mathcal{L}_{\text{matter}}(\psi, \nabla\psi, g, T),
\]
where $R$ is the curvature scalar built from the full (torsion-including)
connection, and the covariant derivative $\nabla$ acts on fermionic fields
$\psi$ with torsion contributions.

\begin{axiom}[There exists a Cartasis density]\label{ax:rhoc}
At sufficiently high mass--energy density, torsion-mediated repulsion exceeds
gravitational attraction. The critical density (the Cartasis density) is
approximately
\[
  \rhoc \sim \frac{m_P^4}{\hbar^3}\,\alpha_{EC},
\]
where $m_P$ is the Planck mass and $\alpha_{EC}$ is a dimensionless
Einstein--Cartan coupling.
\end{axiom}

\noindent Estimates place this at
$\rhoc \sim 10^{50}\,\mathrm{kg\,m^{-3}}$, far below Planck density
($10^{96}\,\mathrm{kg\,m^{-3}}$) but well above nuclear density
($10^{17}\,\mathrm{kg\,m^{-3}}$). This is calculable from
Axioms~\ref{ax:qm}--\ref{ax:ec} in principle; Axiom~\ref{ax:rhoc} is the
explicit acknowledgment that this density exists and matters.

\section{What we do not assume}

\begin{itemize}
  \item No Big Bang as a singular event.
  \item No initial conditions for the universe.
  \item No fine-tuned cosmological constant or inflaton field.
  \item No special low-entropy past as a separate postulate.
  \item No additional dimensions beyond four.
  \item No additional forces or fields beyond Standard Model $+$ gravity $+$ torsion.
  \item No anthropic principle as a separate axiom (it emerges from the framework).
\end{itemize}

\section{What follows automatically}

From Axioms~\ref{ax:qm}--\ref{ax:rhoc} plus infinity (infinite spacetime
volume), the following are consequences rather than assumptions:

\begin{enumerate}
  \item \textbf{Spontaneous OG bounce nucleation.} Rare vacuum fluctuations reach
    Cartasis density. By Axiom~\ref{ax:qm} plus infinity, this happens somewhere
    with probability one. By Axiom~\ref{ax:rhoc}, such fluctuations bounce rather
    than collapsing to singularities.
  \item \textbf{Baby universe production.} The bounce geometry produces new
    spacetime regions (developed in Chapter~\ref{chap:baby}).
  \item \textbf{CPT-symmetric supraverse.} Vacuum fluctuations are statistically
    symmetric in matter/antimatter content. Specific fluctuations have specific
    biases, but the supraverse-wide distribution is symmetric.
  \item \textbf{Thermodynamic equilibrium configuration.} Gravitational entropy
    considerations (high entropy $=$ clumped/black-hole-rich; low entropy $=$
    smooth/uniform) drive the supraverse toward a foam of universes as its
    equilibrium configuration.
  \item \textbf{Local arrow of time.} Each universe has a low-entropy bounce in
    its past and a high-entropy future, producing a local thermodynamic
    time-arrow.
  \item \textbf{Observer placement statistics.} Observers exist in viable
    universes near the peak of the supraverse population distribution.
  \item \textbf{Information preservation.} Hawking radiation entangles parent
    exterior with baby interior. The full supraverse is unitary even though
    individual universes appear to lose information.
\end{enumerate}

\section{On the role of infinity}

The framework requires the supraverse to be effectively infinite---at minimum,
large enough that rare fluctuations to Cartasis density occur. ``Effectively
infinite'' in this context means: large enough relative to the characteristic
timescales of vacuum fluctuation that the relevant fluctuations are
statistically guaranteed to occur.

For a finite supraverse, some fluctuations might never occur and the framework
would have to specify which ones do. This introduces an arbitrary cutoff that we
want to avoid. The natural assumption is true infinity, which removes this
arbitrariness at the cost of being philosophically heavier.

The supraverse being infinite does not require it to be temporally infinite in
some absolute sense---it requires only that enough ``supraverse time'' exist for
the relevant condensation to occur. Whether the supraverse has a beginning at
the supraverse level is a question we do not currently need to answer; we only
need it to be old enough that condensation has reached equilibrium.

\section{On the role of Einstein--Cartan}

Einstein--Cartan is the minimal natural extension of general relativity. It is
what you get by dropping the artificial restriction that the connection must be
symmetric. The connection's antisymmetric part (torsion) is mathematically
natural and couples to a physical quantity that exists in nature (fermion spin).
At ordinary densities, torsion effects are utterly negligible
($\sim 10^{-40}$ corrections to GR). At Cartasis density, they dominate.

We do not postulate Einstein--Cartan as a new theory; we recognize it as the
geometric framework that should have been used all along, with standard general
relativity as the low-density limit. This is not an additional assumption beyond
standard physics---it is the removal of an assumption (the symmetry of the
connection).
